\documentclass[12pt, letterpaper]{article}

% ----------------------------------------------------------------------------------
% PACKAGES
% ----------------------------------------------------------------------------------
\usepackage[margin=0.5in, top=0.5in, bottom=0.5in]{geometry} % Compact Margins
\usepackage[T1]{fontenc}
\usepackage{mathptmx}           % Times New Roman style font
\usepackage{titlesec}
\usepackage{enumitem}
\usepackage{hyperref}
\usepackage{xcolor}

% --- THE MAGIC PACKAGE ---
% ltablex keeps the "tabularx" syntax you like but allows tables to break across pages.
\usepackage{ltablex} 
\keepXColumns % Ensures the X column width is calculated correctly
\setlength{\LTpre}{0pt} % Removes default extra vertical space before tables
\setlength{\LTpost}{0pt} % Removes default extra vertical space after tables

% ----------------------------------------------------------------------------------
% FORMATTING
% ----------------------------------------------------------------------------------
\definecolor{linkblue}{RGB}{32, 64, 151}
\hypersetup{colorlinks=true, linkcolor=linkblue, urlcolor=linkblue}

% Headers: Compact spacing
\titleformat{\section}{\large\bfseries}{}{0em}{}[\vspace{-2pt}\hrule height 1pt \vspace{3pt}]
\titlespacing*{\section}{0pt}{8pt}{4pt}

% Column width for Dates
\newlength{\datecolwidth}
\setlength{\datecolwidth}{2.5cm}

\begin{document}

% ----------------------------------------------------------------------------------
% HEADER
% ----------------------------------------------------------------------------------
\noindent
\begin{minipage}[t]{0.65\textwidth}
    {\huge \textbf{Wanli Qian}} \\
    \vspace{2pt}
    \href{mailto:wanliqia@usc.edu}{wanliqia@usc.edu} $\cdot$ (678)-956-4378 $\cdot$ Los Angeles, CA
\end{minipage}%
\begin{minipage}[t]{0.35\textwidth}
    \raggedleft
    \href{https://www.linkedin.com/in/wanli-qian-7bb399192/}{linkedin} \\
    \href{https://github.com/silverwings-zero}{github} \\
    \href{https://scholar.google.com/citations?user=azFoAIkAAAAJ&hl=en}{Google Scholar Profile}
\end{minipage}

\vspace{0.2cm}

% ----------------------------------------------------------------------------------
% CURRENT APPOINTMENT
% ----------------------------------------------------------------------------------
\section*{Current Appointment}
\noindent
\begin{tabularx}{\textwidth}{@{}p{\datecolwidth} X@{}}
    \textbf{2024--Present} & \textbf{USC HaRVI Lab, PhD Researcher} \\
    & \begin{itemize}[leftmargin=*, nosep, after=\vspace{4pt}]
        \item Investigate how robots can acquire, represent, and generate rich haptic understanding through multimodal sensing, language, and action.
        \item Develop language-conditioned generative models that synthesize realistic tactile and vibration signals from force, acceleration, and vision-based tactile data.
        \item Study active haptic perception, designing exploratory policies that adapt robot actions to minimize uncertainty over object and surface properties.
        \item Build bilateral teleoperation systems with tactile feedback to enable remote surface perception, manipulation, and data collection.
        \item Curate large-scale multimodal haptic datasets (200+ materials) integrating tapping, sliding, audio, and affective descriptors to support scalable learning.
        \item integration of VR/AR and language-based interaction for haptic and texture authoring
    \end{itemize}
\end{tabularx}


% ----------------------------------------------------------------------------------
% EDUCATION
% ----------------------------------------------------------------------------------
\section*{Education}
\noindent
\begin{tabularx}{\textwidth}{@{}p{\datecolwidth} X@{}}
    \textbf{2024--Present} & \textbf{University of Southern California} \\
    & \textit{Ph.D. in Computer Science} \\[4pt]
    \textbf{2022--2024} & \textbf{University of Chicago} \\
    & \textit{Predoctoral M.S. in Computer Science} \\[4pt]
    \textbf{2018--2022} & \textbf{Georgia Institute of Technology} \\
    & \textit{B.S. in Computer Science} \\
\end{tabularx}

% ----------------------------------------------------------------------------------
% PUBLICATIONS
% ----------------------------------------------------------------------------------
\section*{Publications}

\noindent
\begin{tabularx}{\textwidth}{@{}p{\datecolwidth} X@{}}
    
    % --- 2026 ---
    \textbf{2026} & \textbf{Wanli Qian}, Michael Gu, Aiden Chang, Shihan Lu, Heather Culbertson. "Language-Guided Multimodal Texture Authoring via Generative Models." \textit{Proc. IEEE Haptics Symposium (HAPTICS)}, 2026. \textbf{(Distinguished Technical Long Paper Award)}\\
    & \href{https://arxiv.org/abs/2604.06489}{[PDF]} / \href{https://youtu.be/f5vz52uPw-4?si=IwPee25Zb0pjWJ3Y}{[Video]} \\[6pt]

    % --- 2025 ---
    \textbf{2025} & Matthew Jeung, Anup Sathya, \textbf{Wanli Qian}, Steven Arellano, Luke Jimenez, Ken Nakagaki. "Shape n’Swarm: Hands-on, Shape-aware Generative Authoring with Swarm UI and LLMs." \textit{Proc. ACM User Interface Software and Technology (UIST)}, 2025. \\
    & \href{https://dl.acm.org/doi/10.1145/3746059.3747781}{[PDF]} / \href{https://www.youtube.com/watch?v=5u0M9yL7tyY}{[Video]} \\[6pt]

    & Ran Zhou, Jianru Ding, Chenfeng Gao, \textbf{Wanli Qian}, Benjamin Erickson, Madeline Balaam, Daniel Leithinger, Ken Nakagaki. "Shape-Kit: A Design Toolkit for Crafting On-Body Expressive Haptics." \textit{Proc. ACM CHI Conference on Human Factors in Computing Systems}, 2025. \\
     & \href{https://dl.acm.org/doi/10.1145/3706599.3721280}{[PDF]} / \href{https://www.youtube.com/watch?v=W_rZg8f6lZM}{[Video]} \\[6pt]

    % --- 2024 ---
    \textbf{2024} & \textbf{Wanli Qian}*, Chenfeng Gao*, Anup Sathya, Ryo Suzuki, Ken Nakagaki. "SHAPE-IT: Exploring Text-to-Shape-Display for Generative Shape-Changing Behaviors with LLMs." \textit{Proc. ACM User Interface Software and Technology (UIST)}, 2024. \\
    & \href{https://dl.acm.org/doi/10.1145/3654777.3676348}{[PDF]} / \href{https://www.youtube.com/watch?v=2NxzOBc7AdQ}{[Video]} \\[6pt]

    & Chenfeng Gao*, \textbf{Wanli Qian}*, Richard Liu, Rana Hanocka, Ken Nakagaki. "Towards Multimodal Interaction with AI-Infused Shape-Changing Interfaces." \textit{Adjunct Proc. ACM UIST}, 2024. (*Equal Contribution) \\
    & \href{https://dl.acm.org/doi/10.1145/3672539.3686315}{[PDF]} \\[6pt]

    % --- 2022 ---
    \textbf{2022} & Gerry Chen, Sereym Baek, JD Florez, \textbf{Wanli Qian}, Sang-won Leigh, Seth Hutchinson, Frank Dellaert. "GTGraffiti: Spray Painting Graffiti Art from Human Painting Motions with a Cable Driven Parallel Robot." \textit{Proc. IEEE Intl. Conf. on Robotics and Automation (ICRA)}, 2022. \\
    & \href{https://ieeexplore.ieee.org/document/9812008}{[PDF]}/ \href{https://www.youtube.com/watch?v=R4ySYTNGv6s}{[Video]} \\[6pt]

    & \textbf{Wanli Qian}, Zhe Yan, Zhe Zhu, and Wei Yin. 
    ``Weakly Supervised Part-Based Method for Combined Object Detection in Remote Sensing Imagery.'' 
    \textit{IEEE Journal of Selected Topics in Applied Earth Observations and Remote Sensing (JSTARS)}, 
    vol. 15, pp. 5024--5036, 2022. \\ & \href{https://ieeexplore.ieee.org/document/9789416}{[PDF]} \\[6pt]

    % --- 2021 ---
    \textbf{2021} & Hui Wang, Hao Li, \textbf{Wanli Qian}, Wenhui Diao, Liangjin Zhao, Jinghua Zhang, and Daobing Zhang. 
    ``Dynamic Pseudo-Label Generation for Weakly Supervised Object Detection in Remote Sensing Images.'' 
    \textit{Remote Sensing}, vol. 13, no. 8, p. 1461, 2021. \\ & \href{https://www.mdpi.com/2072-4292/13/8/1461}{[PDF]} \\[6pt]

\end{tabularx}

% ----------------------------------------------------------------------------------
% PAST EXPERIENCE
% ----------------------------------------------------------------------------------
\section*{Past Experience}

\noindent
\begin{tabularx}{\textwidth}{@{}p{\datecolwidth} X@{}}
    \textbf{2022--2024} & \textbf{University of Chicago AxLab, Predoctoral Researcher} \\
    & \begin{itemize}[leftmargin=*, nosep, after=\vspace{4pt}]
        \item Explored how natural language can serve as a high-level interface for authoring and controlling dynamic physical shape displays.
        \item Led the transformation of SHAPE-IT into an interactive system enabling on-the-fly 3D shape generation and reconfiguration from conversational language input.
        \item Developed adaptive LLM-based agents that modify shape-display control scripts in response to ongoing user interaction.
        \item Investigated the role of AI-driven authoring tools in expanding the expressivity and usability of tangible and shape-changing interfaces.
    \end{itemize} \\

    \textbf{2022} & \textbf{Kolmostar, Software Engineer Intern} \\
    & \begin{itemize}[leftmargin=*, nosep, after=\vspace{4pt}]
        \item Studied signal propagation and attenuation of wireless microchips in dense urban environments.
        \item Built simulation pipelines to model urban geometry and quantify shadowing effects on signal strength.
        \item Designed interactive visualization tools to support analysis and deployment planning for large-scale urban sensing systems.
    \end{itemize} \\

    \textbf{2021--2022} & \textbf{Georgia Tech Borg Lab, Lab Assistant} \\
    & \begin{itemize}[leftmargin=*, nosep, after=\vspace{4pt}]
        \item Investigated robotic systems as expressive media by combining large-scale cable robots and articulated manipulators with human artistic input.
        \item Developed motion planning and trajectory optimization methods to produce smooth, human-like robotic drawing and writing behaviors.
        \item Explored computational representations of artistic strokes, including polynomial-based models and deformable virtual brush dynamics.
        \item Integrated motion capture and robotic control to analyze and reproduce human artistic trajectories.
    \end{itemize} \\

    \textbf{2019} & \textbf{Megvii, Robotics Software Solution Intern} \\
    & \begin{itemize}[leftmargin=*, nosep, after=\vspace{4pt}]
        \item Researched sampling-based motion planning algorithms for high-DOF industrial manipulators in constrained environments.
        \item Implemented and benchmarked probabilistic and graph-based planners to study trade-offs between efficiency, robustness, and convergence.
        \item Developed visualization tools to analyze planner behavior and performance across different task geometries.
    \end{itemize}
\end{tabularx}

% ----------------------------------------------------------------------------------
% TEACHING
% ----------------------------------------------------------------------------------
\section*{Teaching}

\noindent
\begin{tabularx}{\textwidth}{@{}p{\datecolwidth} X@{}}
    \textbf{2025} & \textbf{Teaching Assistant for CSCI 445: Introduction to Robotics} \\
    & University of Southern California \\ [4pt]

    \textbf{2023} & \textbf{Teaching Assistant for CMSC 25040: Introduction to Computer Vision} \\
    & University of Chicago \\
\end{tabularx}

% ----------------------------------------------------------------------------------
% SERVICE & OUTREACH
% ----------------------------------------------------------------------------------
\section*{Service \& Outreach}

\noindent
\begin{tabularx}{\textwidth}{@{}p{\datecolwidth} X@{}}
    
    % --- Mentorship Programs ---
    \textbf{Mentorship} & \textbf{\href{https://viterbiundergrad.usc.edu/research/curve/}{USC CURVE Program Mentor}} (Spring 2025, Fall 2025, Spring 2026) \\
    & Mentoring undergraduate students in research. \\[4pt]
    
    & \textbf{\href{https://viterbischool.usc.edu/sure/}{USC Viterbi SURE Program Mentor}} (Summer 2025) \\
    & Summer Research Experience for Undergraduates. \\[6pt]

    % --- Outreach ---
    \textbf{Outreach} & \textbf{\href{https://www.peace4kids.org/}{Peace4Kids}} \\
    & Outreach robotics workshop organizer/mentor. \\[6pt]

    % --- Academic Service ---
    \textbf{Reviewing} & CHI, UIST, IEEE Haptics Symposium, IEEE Transactions on Haptics (ToH). \\[4pt]

    \textbf{Volunteering} & RSS Conference Volunteer. \\
    
\end{tabularx}

\end{document}